\resumenCastellano{
El sistema público de salud chileno clasifica sus hospitales por nivel de
complejidad administrativo, clasificación que puede no reflejar la heterogeneidad
de su actividad clínica real. Este estudio identifica perfiles funcionales
empíricos en 65 hospitales públicos de alta y mediana complejidad a partir de su
casuística de Grupos Relacionados por el Diagnóstico (GRD) del período 2019--2024.

Se procesaron 5{.}771{.}841 egresos GRD. Se construyó una matriz de
$65 \times 35$ variables: 27 indicadores clínico-operativos y ocho componentes
principales que resumen la casuística de capítulos CIE-10, secciones CIE-9-MC y
GRD frecuentes ($72{,}3\,\%$ de varianza acumulada). Se aplicó agrupamiento
jerárquico aglomerativo de Ward en dos niveles: un primer nivel con validación
estadística y un segundo exploratorio sobre el núcleo generalista. El pipeline
incluyó análisis de sensibilidad por bloques de variables, comparación con la
clasificación MINSAL y detección multivariada de hospitales atípicos.

La solución de cuatro grupos obtuvo Silhouette $= 0{,}401$ y mediana de
ARI $= 0{,}958$ bajo submuestreo repetido. Se distinguieron tres hospitales
pediátricos, un núcleo generalista de 54 establecimientos, cinco hospitales con
carga heterogénea y estancia prolongada, y tres de especialización heterogénea.
La concordancia con la clasificación MINSAL fue baja (ARI $= 0{,}112$;
NMI $= 0{,}107$), evidenciando que ambas describen dimensiones distintas de la
red. La evaluación en bloques retenidos mostró menor dispersión intragrupo para
Ward $K=4$, aunque sin superioridad robusta al controlar el número de grupos.
El bloque demográfico-obstétrico concentra el $36{,}3\,\%$ de la distancia total;
su exclusión conserva parcialmente los perfiles (ARI $= 0{,}571$). Se
identificaron nueve hospitales atípicos mediante Isolation Forest, distancia al
centroide y Silhouette individual. La subdivisión exploratoria no alcanzó
estabilidad comparable (ARI mediano $= 0{,}559$, $K_{sub} = 9$) y se incluye
como material suplementario.

Los datos GRD permiten describir una heterogeneidad funcional que complementa la
clasificación administrativa, sin establecer jerarquía entre ambas ni justificar
su uso directo en decisiones de financiamiento o evaluación de desempeño.

\vspace*{0.5cm}
\KeywordsES{aprendizaje no supervisado; clustering jerárquico; grupos relacionados por el diagnóstico (GRD); perfilamiento hospitalario; sistema público de salud; selección de variables.}
}

